\section{Related Work}
\label{sec:related_work}
The journey of wind power forecasting (WPF) has been a continuous evolution of improving how we model the complex dynamics of the atmosphere and its interaction with wind turbine technology. This section provides a concise review of the technological milestones that inform the development of PhyST-Net.

\begin{figure}[!t]
\centering
\rule{0.48\textwidth}{0.16\textwidth}
\caption{Technological Timeline of WPF Era. This timeline highlights the transition from statistical foundations to the current physics-aware deep learning era, identifying key architectural breakthroughs.}
\label{fig:timeline}
\end{figure}

\subsection{Temporal Sequence Forecasting in WPF}
\label{sec:related_temporal}
The early evolution of WPF relied primarily on classical linear statistical models, such as Autoregressive Integrated Moving Average (ARIMA) and seasonal variants (SARIMA). While computationally lightweight, these methods assume signal linearity and stationarity, making them ill-suited to capture the turbulent, non-stationary nature of wind signals. SVR and ensemble methods like GBDTs introduced non-linear capacity but were fundamentally spatial-blind and dependent on hand-crafted features, limiting their ability to process high-dimensional spatio-temporal tensors directly.

The rise of deep learning brought recurrent neural networks (RNNs, LSTMs, GRUs) to the forefront, addressing vanishing gradients but introducing sequential bottlenecks for farm-wide optimization. The Transformer architecture \cite{vaswani2017attention} replaced recursion with self-attention, achieving parallelizable training. Models like Informer \cite{zhou2021informer}, Autoformer \cite{wu2021autoformer}, and FEDformer \cite{zhou2022fedformer} reduced complexity or integrated decomposition, yet they frequently focused on global correlations while ignoring local temporal scale structures. To resolve this, PatchTST \cite{nie2023patchtst} and iTransformer \cite{liu2024itransformer} introduced patch-wise processing to preserve semantic information and improve robustness. However, conventional patching uses hard windowing, creating artificial mathematical discontinuities at the boundaries that can distort gradients. Our proposed AGSP module addresses this by introducing learnable Gaussian soft-windows, ensuring that the transition between patches is physically continuous and numerically stable.

\begin{figure}[!t]
\centering
\rule{0.48\textwidth}{0.16\textwidth}
\caption{Attention Maps in Modern WPF Transformers. This visualization demonstrates how multi-head self-attention captures long-range temporal correlations, forming the baseline for our high-capacity temporal tokens.}
\label{fig:attention_vis}
\end{figure}

\subsection{Spatio-Temporal Graph Neural Networks for Wind Power}
\label{sec:related_spatial}
Because wind turbines in a farm are physically coupled through the fluid medium, spatial modeling is as important as temporal modeling. Graph neural networks (GNNs) \cite{kipf2017semi, velickovic2018graph} provide the mathematical tools to represent these couplings. Early works combined graph convolutional networks (GCNs) with temporal filters for traffic or energy forecasting \cite{yu2017spatio, wu2019graph}. In the WPF domain, researchers initially used static graphs based on geographical distance \cite{khodayar2019spatio}, which ignore the anisotropic nature of wind flow. This led to adaptive graph learning \cite{ren2023data, bentsen2022spatio} and dual-graph networks \cite{li2024sdg}. Yet, these models often entangle different spatial scales, failing to distinguish localized wake effects from provincial weather fronts. We propose R-STMF to explicitly decouple these spatial scales.

\subsection{Kolmogorov-Arnold Networks in Scientific Machine Learning}
\label{sec:related_kan}
The black-box nature of traditional multi-layer perceptrons (MLPs) has often limited their adoption in critical power system operations. Kolmogorov-Arnold networks (KAN) \cite{liu2024kan} offer a radical departure from the MLP by replacing fixed node-based activations with learnable spline-based activations on the edges. Recent literature \cite{yamak2025review, han2024kan4tsf, vaca2024kan, xu2024kolmogorov} highlights the potential of KANs for capturing high-order physical dynamics with fewer parameters than standard architectures. In our K-AMC layer, we pass NWP features through a KAN to generate modulation parameters, providing a more physics-aligned feature injection mechanism that captures the complex curvatures of the turbine power curve.

\subsection{Physics-Informed Deep Learning in Renewable Energy}
\label{sec:related_physics}
Finally, the transition from purely data-driven to physics-aware models is a key trend in scientific ML. Physics-informed neural networks (PINNs) regularize the learning process by embedding physical laws into the loss function \cite{welsh2025physics, potosnak2024severe}. However, architectural constraints are often more robust than soft loss-based penalties. Our work follows the philosophy of structural physics induction, embedding physical boundaries directly into the network architecture using the PI-DCH layer to guarantee grid compliance \cite{Shi2023UltraShort}.
